% Source: https://artofproblemsolving.com/community/c3h2646575_is_this_a_coincidence_right_triangle
% Stable identity: aops-topic-2646575-q01
Let the right angle be at $C$, with $BC=a$ and $AC=b$, and let the square be $PQRS$ with $P,Q$ on $AB$ (in the order $A,P,Q,B$), $S$ on $AC$ and $R$ on $BC$. Triangles $APS$ and $BQR$ are right-angled at $P$ and $Q$, with angles $A$ and $B$, and $\tan A=\dfrac ab$, $\tan B=\dfrac ba$. So
\[AP=\dfrac{s}{\tan A}=\dfrac{bs}{a},\qquad QB=\dfrac{s}{\tan B}=\dfrac{as}{b}.\]
Then
\[c=AP+PQ+QB=s\left(\dfrac ba+1+\dfrac ab\right)=s\cdot\dfrac{a^2+b^2+ab}{ab}=s\cdot\dfrac{c^2+ab}{ab},\]
so $\boxed{s=\dfrac{abc}{c^2+ab}}$. For example, for the $3$-$4$-$5$ triangle, $s=\dfrac{60}{37}$ (sketch). (A square in the corner at $C$ instead has side $\dfrac{ab}{a+b}$, which is always larger.)

\begin{asy}
usepackage("amsmath");
size(4.5cm);
defaultpen(fontsize(10pt));
pair C=(0,0),A=(0,3),B=(4,0);real s=60/37;pair u=(B-A)/5,n=(-0.6,-0.8);
pair P=A+s*3/4*u,Q=P+s*u,Sq=P+s*n,R=Q+s*n;
draw(A--B--C--cycle,linewidth(0.8));fill(P--Q--R--Sq--cycle,lightgray);draw(P--Q--R--Sq--cycle);
draw((0.25,0)--(0.25,0.25)--(0,0.25));
label("$A$",A,N);label("$B$",B,E);label("$C$",C,SW);label("$P$",P,NE);label("$Q$",Q,NE);label("$S$",Sq,W);label("$R$",R,S);
label("$b$",(A+C)/2+(0,0.6),W);label("$a$",(B+C)/2+(0.6,0),S);label("$s$",(P+Q)/2,NE,fontsize(8pt));
\end{asy}
